% Generated draft: requires mathematical and source review.
\documentclass[11pt]{article}
\usepackage[margin=25mm]{geometry}
\usepackage{fontspec}
\setmainfont{Latin Modern Roman}
\usepackage{amsmath,amssymb,mathtools,unicode-math}
\setmathfont{Latin Modern Math}
\usepackage{enumitem,needspace,xurl,hyperref}
\setlength{\parindent}{0pt}
\setlength{\parskip}{5pt}
\setlength{\emergencystretch}{3em}
\newcommand{\N}{\mathbb{N}}
\newcommand{\Z}{\mathbb{Z}}
\newcommand{\Q}{\mathbb{Q}}
\newcommand{\R}{\mathbb{R}}
\newcommand{\C}{\mathbb{C}}
\newcommand{\F}{\mathbb{F}}
\newcommand{\A}{\mathbb{A}}
\newcommand{\PP}{\mathbb{P}}
\newcommand{\E}{\mathbb{E}}
\newcommand{\eps}{\varepsilon}
\newcommand{\dd}{\,\mathrm{d}}
\begin{document}
\section*{UnsolvedMath ID: 20000447}
\textbf{Source problem code:} catalogue entry
\subsection{Definitions and mathematical statement}
For an integral coefficient vector $a=(a_1,a_2,a_3,a_4)$, not identically zero, put $A=\max_i|a_i|$. When $\sum_i a_ix_i^3=0$ has a nonzero integral solution, define
\[m(a)=\min\{\max_i|x_i|:0\ne x\in\Z^4,\ \sum_i a_ix_i^3=0\}.\]
One may restrict to primitive $x$ without changing the minimum. The original does not impose coefficient primitivity, cubefreeness or nonzero individual coefficients.

\textbf{Original questions (Swinnerton-Dyer).} Give an upper bound for $m(a)$ in terms of the coefficients whenever the equation is soluble. For $|a_i|<A$, tabulate the least solution sizes and conjecture their growth with $A$. The source compares proposed bounds $m(a)\ll A^{4/3}$ (Swinnerton-Dyer) and $m(a)\ll_\eps A^{1+\eps}$ (Wooley), while Stoll--Stein computations through $|a_i|\le60$ reportedly suggest an $A^2$ upper scale. No modern bound or heuristic is established by these historical statements. Poonen asks for a heuristic, and Swinnerton-Dyer answers no. The ternary quadratic comparison is said to have scale $A$.

The remarks say that, assuming Schinzel's conjecture, finiteness of the relevant Tate--Shafarevich groups, and one unproved lemma, Stoll can produce soluble examples with $m(a)\gg_\eps A^{2-\eps}$; these hypotheses are needed in particular for rational solubility. Their shape is
\[p(x_1^3+2x_2^3)+q(x_3^3+5x_4^3)=0,\]
with primes $p,q\equiv1\pmod3$, $2p$ and $5q$ approximately equal, $2$ a cube modulo $p$ but not $q$, and $5$ a cube modulo $q$ but not $p$, plus unspecified technical conditions. The displayed residue conditions force distinct primes; the omitted conditions and conjectural solubility must not be dropped from the lower-growth claim.

\textbf{Attached unreviewed matching reduction.} For distinct primes $p,q$ with $2$ a noncube modulo $q$ and $5$ a noncube modulo $p$, the aid claims every solution has $x=(qu,qv,pw,pz)$ and satisfies, for a common integer $t$,
\[u^3+2v^3=p^2t,\qquad w^3+5z^3=-q^2t.\]
For $B\ge0$ define the finite integer sets
\[T_1(B)=\left\{\frac{u^3+2v^3}{p^2}\in\Z:|u|,|v|\le\lfloor B/q\rfloor\right\},\quad T_2(B)=\left\{-\frac{w^3+5z^3}{q^2}\in\Z:|w|,|z|\le\lfloor B/p\rfloor\right\}.\]
The reported exact existence criterion is $(T_1(B)\cap T_2(B))\setminus\{0\}\ne\varnothing$ for a nonzero solution of height at most $B$; primitive existence at that bound is equivalent by dividing a common factor. The square-modulus congruences also give the reported lower-bound framework through their nonzero solution minima. Distinctness is explicit here because the reduction modulo one prime needs the other coefficient invertible; the aid short summary does not state it separately.

The aid gives the exact claimed certification
\[m(31,62,13,65)=248,\qquad (169,91,-248,31)\text{ attains it primitively}.\]
It says complete bounded matching tables prove minimality, but those tables are not in the supplied record. This one finite example and the matching reduction do not establish any of the proposed general growth exponents. All aid assertions require independent review.
\subsection{Short English statement}
The source asks for general least-solution bounds and coefficient-height data for diagonal cubic surfaces, retaining competing historical exponents and conditional lower-growth examples. The aid supplies a square-modulus matching criterion for a restricted prime family and reports one exactly certified least height, without resolving the general growth question.
\subsection{Sources}
\begin{itemize}
\item Local UnsolvedMath registry, permanent ID 20000447, source code catalogue entry; \texttt{aim-algebraic-number-theory-notes.json}, notes[98], tag \texttt{question}.
\item AIM workshop \emph{Rational and integral points on higher dimensional varieties}, Question 14 (Swinnerton-Dyer): \url{https://aimath.org/WWN/qptsurface2/qptsurface2.pdf}. Original remarks attribute proposed exponents to Swinnerton-Dyer and Wooley, computations to Stoll and Stein, conditional examples to Stoll, and a heuristic question to Poonen. No publications or technical lemma are identified.
\item Attached machine-generated research attempt 1: result \texttt{partial\_result}, problem status at run \texttt{open}, rigor label \texttt{proved}, novelty confidence low. It is not peer reviewed and requires independent verification before citation; the exact matching tables and external references are not supplied. No PDF was consulted.
\end{itemize}
\end{document}
